\section{Pooled Class Results}
\label{sec:pooled}

\subsection{Choice Probabilities (3a)}
The pooled data contain 30,450 choices, with $n_j$ right choices and $m_j$
left choices in bin $j$. Treating the supplied mean $d_j$ as each trial's
value difference gives the approximate log likelihood
\begin{equation}\label{eq:pooled-likelihood}
  LL_c(\alpha)
  \;=\;\sum_{j=1}^{15}
  \left[n_j\log P_\alpha(d_j)+m_j\log P_\alpha(-d_j)\right].
\end{equation}
Binomial coefficients can be omitted because they do not depend on $\alpha$.
The subscript $c$ denotes the class data. Maximizing
equation~\eqref{eq:pooled-likelihood} gives $\widehat\alpha_c=\AlphaPooled$, with
$LL_c(\widehat\alpha_c)=\LLPooled$.

The plotted frequencies are $p_j=n_j/(n_j+m_j)$ and
$e_j=P_{\widehat\alpha_c}(d_j)$. Under the bin-mean assumption,
$P_{\widehat\alpha_c}^{-1}(e_j)=d_j$ exactly, so the supplied means are also
the inverse-probability coordinates used in Question 1(d).

\begin{figure}[!ht]
  \centering
  \caption{\capttitle{Pooled Choice Frequencies}}
  \label{fig:pooled-choice}
  \includegraphics[width=0.94\textwidth]{fig-pooled-choice.pdf}
  \floatnotes{Source: \texttt{popData.csv}; 30,450 choices in 15 supplied bins.
  Open circles show $n_j/(n_j+m_j)$; black dots and the curve show
  predictions using the class estimate of alpha.
  Estimation treats each bin's supplied mean value difference as the value
  difference on every trial in that bin.}
\end{figure}

\begingroup
\clubpenalty=10000
The class choice fit has a count-weighted RMSE of \ChoiceRMSEPooled\ percentage
points (Figure~\ref{fig:pooled-choice}). The errors are systematic.
Near $d_j=-1.35$ and $d_j=1.29$, observed right-choice shares
are 26.5 and 75.6 percent; the fitted shares, 30.4 and 69.0 percent, lie
closer to one-half. This pattern reverses in the tails. Even with a value
advantage of \$7.81, the right snack was passed over in 5.3 percent of choices,
against just 0.8 percent predicted. A single $\alpha$ cannot match both the
steep change around equal bids and the less extreme choices in the tails.
\par
\endgroup

\subsection{Mean Response Times (3b)}
The pooled file contains only bin means, so individual outliers cannot be trimmed.
With $\widehat\alpha_c$ fixed, I regress $T_j$ on
$g_j=G_{\widehat\alpha_c}(d_j)$, giving each of the 15 bins equal weight.
The estimates are
\[
  \widehat A_c=\TimeScalePooled,
  \qquad \widehat t_{0,c}=\OffsetPooled\ \text{ms},
  \qquad R^2=\RSquaredPooled.
\]

\begin{figure}[!ht]
  \centering
  \caption{\capttitle{Pooled Mean Response Times}}
  \label{fig:pooled-time}
  \includegraphics[width=0.94\textwidth]{fig-pooled-response-time.pdf}
  \floatnotes{Source: \texttt{popData.csv}; 15 supplied response-time means,
  in milliseconds. Open circles show $T_j$; black dots show
  $\widehat A_c g_j+\widehat t_{0,c}$.
  The time function uses the class estimate of alpha, and all bins receive
  equal OLS weight. Under the imposed bin-mean
  approximation, the black dots lie exactly on the theoretical curve.}
\end{figure}

The class means closely follow the predicted inverted-U
(Figure~\ref{fig:pooled-time}). Their RMSE is
\RTRMSEPooled\ milliseconds, compared with \RTRMSEIndividual\ milliseconds
for my data; $R^2$ is 0.961 (Table~\ref{tab:estimates}).

The largest residual is in the bin near $d_j=-0.03$, whose mean of
1,438 milliseconds exceeds the fitted value by about 80 milliseconds.
At the far right, the class also takes about 74 milliseconds longer
than predicted.

\begin{table}[!ht]
  \centering
  \caption{\capttitle{Estimated Parameters and Descriptive Fit}}
  \label{tab:estimates}
  \small
  \begin{tabular}{@{}lrr@{}}
    \toprule
    & Individual & Pooled class \\
    \midrule
    Choice trials &       435 &    30,450 \\
Choice sensitivity & 0.408351 & 0.308202 \\
Log likelihood & \(    -192.608\) & \(-14{,}580.856\) \\
Choice RMSE (percentage points) &   6.70 &   3.64 \\
\midrule
Response-time bins & 15 & 15 \\
Time scale, $\widehat{A}$ &  2,913.033 &  2,656.594 \\
Nondecision time, $\widehat{t}_0$ (ms) &    644.123 &    539.137 \\
\(R^2\) of bin means & 0.5542 & 0.9607 \\
Response-time RMSE (ms) &   172.21 &    36.71 \\
\bottomrule

  \end{tabular}
  \floatnotes{Choice estimates $\widehat\alpha$ (individual) and
  $\widehat\alpha_c$ (class) use maximum likelihood with no intercept.
  Response-time estimates use unweighted OLS on 15 means, holding the
  relevant choice estimate fixed. Individual response-time means use
  \RTKept\ retained trials; pooled means are used as supplied.
  Choice RMSE is $[\sum_j N_j(p_j-e_j)^2/\sum_jN_j]^{1/2}$,
  with $N_j$ the total choice count in a bin. Time RMSE divides the
  residual sum of squares by 15. $\alpha$ is in inverse dollars and
  $A$ is in milliseconds times dollars.}
\end{table}
